\documentclass[../../../../main.tex]{subfiles}
\begin{document}

\begin{flushright}
\footnotesize\textit{【自動文字起こし・要確認】}
\end{flushright}

[解]\\
(i)
\[
  p = \sum_{k=1}^6 p_k^2
\]

(ii)
\begin{align}
  p &= p_1^2 + p_2^2 + p_3^2 + p_4^2 + p_5^2 + p_6^2 \label{eq4_1} \\
  1 &= p_1 + p_2 + p_3 + p_4 + p_5 + p_6 \label{eq4_2}
\end{align}

コーシー・シュワルツの不等式から
\[
  (p_1+p_2+p_3+p_4+p_5+p_6)^2 \le 6p
\]
\[
  \therefore p \ge \frac{1}{6} \ (\because \text{\eqref{eq4_2}}) \quad \text{終}
\]
等号成立は$p_1 = p_2 = p_3 = p_4 = p_5 = p_6$の時で, \eqref{eq4_2}からこの時
\[
  p_k = \frac{1}{6}
\]
である 終

\end{document}
